\documentclass{article}

\usepackage{booktabs}
\usepackage{tabularx}
\usepackage{array}
\usepackage{longtable}
\usepackage{pdflscape}
\usepackage[round]{natbib}
\usepackage{tikz}
\usetikzlibrary{positioning, arrows.meta}
\usepackage{hyperref}

\hypersetup{
    colorlinks=true,       % false: boxed links; true: colored links
    linkcolor=red,          % color of internal links (change box color with linkbordercolor)
    citecolor=green,        % color of links to bibliography
    filecolor=magenta,      % color of file links
    urlcolor=cyan           % color of external links
}

\title{Hazard Analysis\\\progname}

\author{\authname}

\date{}

\input{../Comments.text}
\input{../Common.text}

\begin{document}

\maketitle
\thispagestyle{empty}

~\newpage

\pagenumbering{roman}

\begin{table}[hp]
\caption{Revision History} \label{TblRevisionHistory}
\begin{tabularx}{\textwidth}{llX}
\toprule
\textbf{Date} & \textbf{Developer(s)} & \textbf{Change}\\
\midrule
14 October 2026 & Rayan Nasrallah & Added critical assumptions CA-001 to
CA-006 and hazards H-RAY-001 to H-RAY-007 (\#91).\\
15 October 2026 & Rayan Nasrallah & Added my reflection responses to the appendix and my generative AI use disclosure (\#140).\\
15 October 2026 & Ahmed Elzaria & Added drawing-input and guided-lesson FMEA hazards H-AHM-001 to H-AHM-007 (part of \#135).\\
16 October 2026 & Muhammad Huzaifah & Listed the HA-derived requirements with their hazards in Section 6 and planned their phases in the Section 7 roadmap (\#128).\\
\bottomrule
\end{tabularx}
\end{table}

~\newpage

\tableofcontents

~\newpage

\pagenumbering{arabic}

\wss{You are free to modify this template.}

\section{Introduction}

\wss{You can include your definition of what a hazard is here.}

\section{Scope and Purpose of Hazard Analysis}

\wss{You should say what \textbf{loss} could be incurred because of the
hazards.}

\section{System Boundaries and Components}

\wss{Dividing the system into components will help you brainstorm the hazards.
You shouldn't do a full design of the components, just get a feel for the major
ones.  For projects that involve hardware, the components will typically include
each individual piece of hardware.  If your software will have a database, or an
important library, these are also potential components.}

\section{Critical Assumptions}

The assumptions below are kept few and are written so that they do not
remove hazards. Squigglr does not assume that an adult watches the child
during practice, that the network is always available, or that the device
meets Squigglr's requirements; those conditions are treated as failure
modes in the Failure Mode and Effect Analysis (FMEA). Sources are cited by
their identifiers in the Software Requirements Specification (SRS),
Appendix A, and interview turns by speaker and timestamp, for example
(SRC-005, P1, 34:13). Each assumption lists the H-RAY hazards it affects.

\begin{description}
\item[CA-001] An adult sets up Squigglr for a child, chooses and starts
the lesson and its starting stage, and is the only person meant to accept
a suggestion to go back a stage. Within a lesson the child does not
navigate, exit or restart; Squigglr moves between stages. An adult may not
be present or watching while the child practises.\\
\emph{Rationale:} children should practise without constant adult help,
while adults handle setup (SRC-005, P1, 34:13), and adults decide progression
(SRC-005, P1, 28:37; CON-007).\\
\emph{Hazards affected:} H-RAY-001, H-RAY-004, H-RAY-005.

\item[CA-002] The device, browser, touchscreen or stylus in use may not be
one of the supported combinations (\texttt{SUPPORTED\_DEVICES}, TBD), and
Squigglr may be opened on a device without a stylus, where the child can
draw with a finger but not as intended (ASM-003).\\
\emph{Rationale:} Squigglr is a website opened from a web address with no
installation (CON-001, CON-002), so the team does not control the device.
The supported combinations are not yet defined (CON-003), and the team did
not promise that Squigglr works on every device (SRC-005, Moustafa,
6:03).\\
\emph{Hazards affected:} H-RAY-002.

\item[CA-003] The network connection may be slow or lost at any time,
including during an attempt.\\
\emph{Rationale:} attempts are saved to an account and reviewed later from
another device, and offline access was not promised (SRC-005, Moustafa,
6:03). Home and school networks are outside the team's control.\\
\emph{Hazards affected:} H-RAY-002, H-RAY-006.

\item[CA-004] One device may be used by several children, at school or at
home.\\
\emph{Rationale:} school devices may be shared by several children
(SRC-004), and Squigglr may be used at home, at school or in other settings
such as rehabilitation centres (SRC-005, P2, 21:25).\\
\emph{Hazards affected:} H-RAY-006.

\item[CA-005] The adults who read practice records are not assumed to have
training in visual-motor development, and they may show the records to
other people, such as teachers or therapists.\\
\emph{Rationale:} parents and caregivers are not assumed to be specialists
(SRS, Hands-On Users), and records exist to support conversations about
what a child needs (SRC-005, Moustafa, 6:03). Squigglr makes no diagnosis
or treatment claims (SRC-001).\\
\emph{Hazards affected:} H-RAY-003.

\item[CA-006] The supporting adult chooses the adult-set values, such as
\texttt{STAGE\_PASS}, \texttt{DIFFICULTY\_STREAK} and
\texttt{ATTEMPT\_LIMIT}, for each child within a range the team sets
(SRC-008). The ranges are not found from child testing, and the value an
adult picks, or the default, may not suit their child.\\
\emph{Rationale:} no information is gathered from children and no study
with children is run (CON-005, CON-006), so the ranges cannot be tuned on
children (SRS, Goals of the Project). Children reach these skills at
different points (SRC-005, P1, 20:27).\\
\emph{Hazards affected:} H-RAY-004, H-RAY-005, H-RAY-007.
\end{description}

\section{Failure Mode and Effect Analysis}

\wss{Include your FMEA table here. This is the most important part of this document.}
\wss{The safety requirements in the table do not have to have the prefix SR.
The most important thing is to show traceability to your SRS. You might trace to
requirements you have already written, or you might need to add new
requirements.}
\wss{If no safety requirement can be devised, other mitigation strategies can be
entered in the table, including strategies involving providing additional
documentation, and/or test cases.}

\subsection{Supervision, Context of Use and Use of Records (H-RAY)}

Table~\ref{tab:fmea-ray} covers hazards from how Squigglr is supervised,
where and on what it is used, and how its records are read. Each
recommended action starts with Eliminate, Mitigate or Document/Test.
Mitigating requirements trace to the goals (G.\#) and constraints
(CON-\#\#\#) in the SRS, or to safety requirements SAF-RAY-\#\#\# in the SRS,
Safety-Critical Requirements.

\begin{landscape}
{\footnotesize
\setlength{\tabcolsep}{3pt}
\begin{longtable}{@{}
  >{\raggedright\arraybackslash}p{1.4cm}
  >{\raggedright\arraybackslash}p{1.8cm}
  >{\raggedright\arraybackslash}p{2.4cm}
  >{\raggedright\arraybackslash}p{2.3cm}
  >{\raggedright\arraybackslash}p{2.6cm}
  >{\raggedright\arraybackslash}p{1.9cm}
  >{\raggedright\arraybackslash}p{3.4cm}
  >{\raggedright\arraybackslash}p{2.0cm}@{}}
\caption{FMEA: supervision, context of use and use of records}
\label{tab:fmea-ray}\\
\toprule
\textbf{Hazard ID} & \textbf{Component} & \textbf{Failure mode} &
\textbf{Cause(s)} & \textbf{Effect and loss} & \textbf{Category} &
\textbf{Recommended action} & \textbf{Mitigating requirement IDs} \\
\midrule
\endfirsthead
\toprule
\textbf{Hazard ID} & \textbf{Component} & \textbf{Failure mode} &
\textbf{Cause(s)} & \textbf{Effect and loss} & \textbf{Category} &
\textbf{Recommended action} & \textbf{Mitigating requirement IDs} \\
\midrule
\endhead
\bottomrule
\endlastfoot
H-RAY-001 & Lesson flow; adult controls & A child practises with no
adult watching, or the adult is inattentive, and the child reaches an
adult-only action. & Children are meant to practise without constant adult
help (CA-001); adult controls use the same touch interface as the child. &
The child accepts a suggestion to go back a stage, or changes lesson or
stage, without an adult's decision; the child's distress goes unnoticed.
Loss: progression not chosen by an adult; well-being. &
Psychological/ well-being &
Mitigate: place every adult-only action (choosing a lesson or stage,
accepting a suggestion) behind an adult check that a child cannot pass by
tapping icons; keep a pending suggestion waiting, without blocking
practice, until an adult reviews it. &
CON-007, G.3, SAF-RAY-001 \\
\addlinespace
H-RAY-002 & Browser front end; saving of attempts & Squigglr is used on an
unsupported device or browser, without a stylus, or loses its connection
during an attempt, at home, at school or on a shared device. & Device not
controlled by the team (CA-002); no offline access (CA-003). & Strokes
are not captured correctly, the child's attempt is lost, or a lesson
cannot be completed; the child is frustrated and the record is incomplete.
Loss: lost practice data; availability. &
Data integrity or Availability &
Mitigate: before a lesson starts, check the device and browser against
\texttt{SUPPORTED\_DEVICES} and show the adult a notice when they do not
match; keep each attempt until the server confirms it is saved, and tell
the adult when an attempt could not be saved. &
CON-001, CON-002, CON-003, G.5, SAF-RAY-002, SAF-RAY-003 \\
\addlinespace
H-RAY-003 & Attempt records; progress summaries & Practice records are
used as clinical or diagnostic evidence, or as a measure of the child's
ability. & Adults reading records lack specialist training, and records
may be passed to others (CA-005); fewer attempts may be read as better
skill (SRC-005, Moustafa, 6:03). & An adult or professional draws a wrong
conclusion about the child, such as a label or a placement decision.
Loss: harm to the child's well-being and to decisions about support. &
Psychological/ well-being &
Eliminate: show no diagnostic labels, judgements of the child's abilities
or comparisons with other children in any record or summary. Mitigate: state on every summary and
export that the records describe practice only and are not an assessment.
&
G.5, G.6, SAF-RAY-004 \\
\addlinespace
H-RAY-004 & Lesson flow; feedback & Repeated failed attempts at one stage
continue until the child becomes frustrated or agitated. & Failure
feedback repeats; \texttt{ATTEMPT\_LIMIT} may be too high for a given
child (CA-006); an improvement extension lengthens the lesson; no adult may be watching (CA-001). & The child becomes
distressed, stops practising or refuses the activity later. Loss:
well-being; the practice benefit. &
Psychological/ well-being &
Mitigate: end the lesson at \texttt{ATTEMPT\_LIMIT}, or at
\texttt{ATTEMPT\_LIMIT\_MAX} when extended, and save progress (G.3); let audio feedback be muted (SRC-005, P2, 37:39). Document/Test: test that each lesson ends at
\texttt{ATTEMPT\_LIMIT} and that no extended lesson passes
\texttt{ATTEMPT\_LIMIT\_MAX}. &
G.3, SAF-RAY-005 \\
\addlinespace
H-RAY-005 & Lesson flow; attempt records & After a lesson ends at
\texttt{ATTEMPT\_LIMIT}, it is restarted at once and the repetition
continues. & No adult may be watching (CA-001); a restart starts a new
count of attempts. & The limit does not stop the
repetition, and the adult cannot see how often it happened. Loss:
well-being; misleading records. &
Psychological/ well-being &
Mitigate: record each lesson that ended at \texttt{ATTEMPT\_LIMIT} and
show the adult how many times this happened per lesson; keep earlier
attempts when a lesson is restarted. Document/Test: test that a restart
after the limit keeps the earlier record. &
G.3, G.5, SAF-RAY-006 \\
\addlinespace
H-RAY-006 & Learner selection; attempt records & On a shared device,
attempts or a pending suggestion are saved to, or shown for, the wrong
child. & Several children use one device (CA-004); the active child is
not changed between sessions; a lost connection leaves a session open
(CA-003). & One child's records are shown to another child or adult, or a
child's history is mixed with another's, so an adult accepts a suggestion
based on the wrong attempts. Loss: privacy; record integrity. &
Privacy/ security &
Mitigate: show which learner is active on every practice screen in a way
a child who cannot read can recognize; require an adult action to switch
learners; show the attempts behind a suggestion when the adult reviews it.
&
G.5, SAF-RAY-001, SAF-RAY-007 \\
\addlinespace
H-RAY-007 & Child profile settings; lesson flow & A stage check or lesson
limit uses a value that does not suit the child: \texttt{STAGE\_PASS} too
strict or too lenient, or \texttt{ATTEMPT\_LIMIT} not larger than
\texttt{DIFFICULTY\_STREAK}, or \texttt{ATTEMPT\_LIMIT\_MAX} below
\texttt{ATTEMPT\_LIMIT}. & The adult misjudges the example drawings or
leaves a default that does not suit the child (CA-006); a value is saved
outside its range; a child changes a setting. & A child cannot pass a stage
they can draw, or passes one they cannot, and the records mislead the
adult. Loss: well-being; misleading records. &
Psychological/ well-being &
Mitigate: accept only values within the team-set range, and only an
\texttt{ATTEMPT\_LIMIT} larger than \texttt{DIFFICULTY\_STREAK}; show
example drawings of each end of the \texttt{STAGE\_PASS} range; place
settings behind the adult check; show the values in use next to the
child's records. Document/Test: test that out-of-range and invalid values
are refused and that the default is the midpoint. &
G.2, G.3, SAF-RAY-008 \\
\end{longtable}
}
\end{landscape}

\subsection{Drawing Input and Guided-Lesson Decisions (H-AHM)}
\label{sec:fmea-ahm}
Table~\ref{tab:fmea-ahm} examines failures in the capture and interpretation
of stylus input and in the decisions that use it. These failures can occur
even on an otherwise supported device with a working network, so they are
distinct from the environmental failure in H-RAY-002. Effects are assessed
against drawing accuracy, record integrity and the child's well-being;
Squigglr does not make clinical assessments (BR-006). Recommended actions
are traced to the learning-activity requirements in the SRS. Two additional
safeguards, SAF-AHM-001 and SAF-AHM-002, are specified in the SRS
Safety-Critical Requirements section.
\begin{landscape}
{\footnotesize
\setlength{\tabcolsep}{2pt}
\begin{longtable}{@{}
  >{\raggedright\arraybackslash}p{1.85cm}
  >{\raggedright\arraybackslash}p{1.8cm}
  >{\raggedright\arraybackslash}p{2.4cm}
  >{\raggedright\arraybackslash}p{2.3cm}
  >{\raggedright\arraybackslash}p{2.6cm}
  >{\raggedright\arraybackslash}p{1.9cm}
  >{\raggedright\arraybackslash}p{3.4cm}
  >{\raggedright\arraybackslash}p{2.1cm}@{}}
\caption{FMEA: drawing input and guided-lesson decisions}
\label{tab:fmea-ahm}\\
\toprule
\textbf{Hazard ID} & \textbf{Component} & \textbf{Failure mode} &
\textbf{Cause(s)} & \textbf{Effect and loss} & \textbf{Category} &
\textbf{Recommended action} & \textbf{Mitigating requirement IDs} \\
\midrule
\endfirsthead
\toprule
\textbf{Hazard ID} & \textbf{Component} & \textbf{Failure mode} &
\textbf{Cause(s)} & \textbf{Effect and loss} & \textbf{Category} &
\textbf{Recommended action} & \textbf{Mitigating requirement IDs} \\
\midrule
\endhead
\bottomrule
\endlastfoot
\mbox{H-AHM-001} & Stylus input capture; attempt state & A stroke or part of
one stroke disappears during a multi-stroke attempt, but the attempt is
still finalized and evaluated. & Missed pointer samples, a brief
input interruption or premature finalization (ASM-004); a stylus lift
is mistaken for the end of an attempt. & The drawing appears incomplete
or is scored inaccurately; a child receives undeserved failure feedback
or loses a practice attempt. Loss: accurate feedback and record integrity. &
Data integrity & Mitigate: retain ordered points and stroke boundaries;
distinguish stylus lifts from finalization; treat interrupted or
unevaluable input as non-passing without consuming an attempt.
Document/Test: replay partial strokes and interruptions before and
during a stroke. & FR-LRN-007, FR-LRN-008, FR-LRN-009, FR-LRN-010 \\
\addlinespace
\mbox{H-AHM-002} & Stylus event capture; attempt evaluation & One drawn stroke
is captured or evaluated twice, inflating coverage or producing a
spurious extra line. & Duplicate input events, accidental replay after
reconnection, or combining the same stroke twice. & A poor attempt may
pass, or a correct one may fail; saved strokes no longer represent what
was drawn. Loss: integrity of feedback and practice records. &
Data integrity & Mitigate: reject duplicate events within the same
attempt before evaluation, without deleting legitimate separate
strokes. Document/Test: replay duplicate events and distinct repeated
strokes; verify that each is handled appropriately. &
FR-LRN-007, FR-LRN-008; SAF-AHM-001 \\
\addlinespace
\mbox{H-AHM-003} & Input coordinates; reference-shape alignment & Valid
stylus points are mapped to the wrong position or order and used to
score the drawing. & Wrong canvas scale, coordinate offset or
incorrect ordering after a resize or orientation change (ASM-004). &
An accurate attempt fails, an inaccurate attempt passes, or incorrect
strokes are saved. Loss: data integrity and trustworthy progress
records. & Data integrity & Mitigate: verify that captured coordinates
remain in the current drawing surface and retain their recorded order;
reject input that cannot be mapped reliably. Document/Test: repeat
known stroke traces before and after resizing or rotating the device. &
FR-LRN-007, FR-LRN-010; SAF-AHM-002 \\
\addlinespace
\mbox{H-AHM-004} & Completeness and pass classifier; lesson flow & An
incomplete drawing (for example, a square missing a side) is reported
as passing and advances the child. & The geometric result is high even
though a required segment is missing, or completeness is not checked
before threshold comparison. & The system advances despite an
unfinished task and records a misleading success. Loss: record
integrity and appropriate guidance. & Data integrity & Eliminate:
require completeness and the adult-selected passing threshold to both
hold before a pass; otherwise remain at the same stage. Document/Test:
use complete and missing-part drawings immediately above and below
the threshold, including a high-scoring incomplete example. &
FR-LRN-011, FR-LRN-013, FR-LRN-023 \\
\addlinespace
\mbox{H-AHM-005} & Feedback-selection logic; stage-result mapping & A detected
attempt result or drawing difficulty is associated with the wrong
pre-approved feedback observation (for example, a missed endpoint is
reported as a missing segment). & A classification identifier is mapped
to the wrong feedback entry, or the feedback lookup uses the wrong lesson
or guidance stage. & The child is directed to correct a part of the
drawing that was not identified as incorrect; the reported observation
does not match the evaluated attempt. Loss: accurate guidance and
feedback integrity. & Data integrity & Mitigate: select feedback from
the approved catalogue using the evaluated result and correct lesson and
stage identifiers; when no matching observation exists, do not present
an unrelated correction. Document/Test: check the mapped message for
every predefined classification and each drawing stage; test missing and
mismatched identifiers. & FR-LRN-011, FR-LRN-013, FR-LRN-014 \\
\addlinespace
\mbox{H-AHM-006} & Failure-streak and suggestion logic; adult decision &
Squigglr recommends moving backward at the wrong time, or changes
the stage without the adult accepting the suggestion. & Invalid input
is counted as failure, the streak is not reset, or the improvement
suppression rule is omitted (BR-003). & The child receives unsuitable
support or a needless interruption; adult control is bypassed. Loss:
well-being and authorized lesson progression. & Psychological/ well-being &
Mitigate: base the suggestion only on counted non-passing attempts and
the defined improvement predicate; leave the current stage unchanged
until an authorized adult accepts. Document/Test: replay improving,
non-improving, invalid and threshold-boundary sequences; decline or
ignore a suggestion. & FR-LRN-009, FR-LRN-018, FR-LRN-019,
FR-LRN-020 \\
\addlinespace
\mbox{H-AHM-007} & Learner settings; attempt and lesson evaluation &
The active attempt is evaluated using another learner's threshold, a
stale value or a setting different from the adult's saved selection. &
Wrong profile linkage, settings not refreshed after an adult change,
or defaults applied over saved choices (BR-008; CA-006). & A drawing
passes or fails incorrectly, a lesson ends too early or continues too
long, and the adult receives misleading records. Loss: well-being and
record integrity. & Data integrity & Mitigate: use the active learner's
validated saved settings for passing, suggestions and limits; use
midpoint defaults only for unset values. Document/Test: switch learner
profiles and change settings; test boundaries, invalid combinations and
the exact saved-versus-default choice. & FR-LRN-016, FR-LRN-017,
FR-LRN-019, FR-LRN-021 \\
\end{longtable}
}
\end{landscape}

\section{Safety and Security Requirements}

This section lists the safety and security requirements derived from the
hazard analysis (HA). A requirement is HA-derived when its ID appears in
the Mitigating requirement IDs column of any of the three per-author FMEA
tables (H-RAY, H-AHM and H-MUH) in the Failure Mode and Effect Analysis
section, so that it carries out the recommended action of at least one
hazard. As the teaching assistant (TA) advised
(SRC-006), each requirement is written once, in the SRS, together with its
rationale, fit criterion, priority,
phase and sources. This section therefore adds no requirement of its own:
Table~\ref{tab:safety-reqs} repeats the statement of each HA-derived
requirement exactly as the SRS states it.

The table has two groups. Functional requirements come first, in the order
of the SRS subsections that hold them: the Learning Activity, Creative
Drawing, Attempt Record and Progress Summary requirements.
Safety-critical requirements follow (SAF-RAY and SAF-AHM); the SRS keeps
them in its Safety-Critical Requirements subsection because
they do not fit another section. For each requirement, the table gives
its ID, with the short title the SRS gives it where there is one, its
statement, the hazards whose FMEA rows cite it, and the SRS subsection that
holds it. The subsection is named in words, because the SRS is a separate
document; for the same reason, a statement that refers to an SRS table
names the table by its caption. Names in monospaced capitals, such as \texttt{STAGE\_PASS}, are
symbolic constants that the SRS defines in its table ``Central symbolic
constants register''.

Some rows of the H-RAY table also cite goals (G.2, G.3, G.5 and G.6) or
constraints (CON-001, CON-002, CON-003 and CON-007) among their
mitigations. A goal states what \progname{} is for, and a constraint is
imposed on the project from outside, so neither is a requirement and
neither is repeated here. Every row that cites a goal or a constraint also
cites at least one safety-critical requirement of
Table~\ref{tab:safety-reqs}.

{\footnotesize
\setlength{\tabcolsep}{4pt}
\exhyphenpenalty=10000 % keep IDs such as FR-CRT-017 on one line
\begin{longtable}{@{}
  >{\raggedright\arraybackslash}p{2.1cm}
  >{\raggedright\arraybackslash}p{\dimexpr\linewidth-5.75cm-6\tabcolsep\relax}
  >{\raggedright\arraybackslash}p{1.75cm}
  >{\raggedright\arraybackslash}p{1.9cm}@{}}
\caption{HA-derived requirements, as stated in the SRS}
\label{tab:safety-reqs}\\
\toprule
\textbf{ID} & \textbf{Requirement} & \textbf{Hazards} &
\textbf{SRS section} \\
\midrule
\endfirsthead
\toprule
\textbf{ID} & \textbf{Requirement} & \textbf{Hazards} &
\textbf{SRS section} \\
\midrule
\endhead
\midrule
\multicolumn{4}{@{}r@{}}{\emph{Continued on the next page}} \\
\endfoot
\bottomrule
\endlastfoot
\multicolumn{4}{@{}l}{\textbf{Functional requirements}} \\*
\addlinespace
FR-LRN-007\newline\emph{Stylus stroke acquisition} & Squigglr shall accept
stylus-drawn strokes as ordered sequences of input points for the active lesson
attempt and preserve their order and stroke boundaries for subsequent
classification. & H-AHM-001, H-AHM-002, H-AHM-003 & Learning Activity \\
\addlinespace
FR-LRN-008\newline\emph{Multi-stroke attempts} & Squigglr shall evaluate the
collected strokes of a single attempt together rather than classify each stylus
lift as a completed attempt. & H-AHM-001, H-AHM-002 & Learning Activity \\
\addlinespace
FR-LRN-009\newline\emph{Interrupted attempt handling} & Squigglr shall
distinguish an interrupted or unfinished drawing attempt from a completed
drawing attempt, shall not record it as passed, and shall not consume an
evaluable attempt from the stage limit. & H-AHM-001, H-AHM-006 & Learning
Activity \\
\addlinespace
FR-LRN-010\newline\emph{Invalid and ambiguous input} & Squigglr shall withhold
passing status from an attempt whose required stylus input is missing, invalid
or too ambiguous to evaluate against the stage criteria, and shall not count
such input toward the stage's evaluable-attempt or unsuccessful-attempt limits.
& H-AHM-001, H-AHM-003 & Learning Activity \\
\addlinespace
FR-LRN-011\newline\emph{Completeness classification} & Squigglr shall classify
every evaluable drawing attempt as complete or incomplete relative to the
selected lesson and guidance stage. & H-AHM-004, H-AHM-005 & Learning Activity
\\
\addlinespace
FR-LRN-013\newline\emph{Pass classification} & Squigglr shall report an
evaluable attempt as passing its stage only when the drawing is classified as
complete and its result meets the current learner's \texttt{STAGE\_\allowbreak
PASS} criterion. & H-AHM-004, H-AHM-005 & Learning Activity \\
\addlinespace
FR-LRN-014\newline\emph{Predefined difficulty feedback} & Squigglr shall present
only pre-approved, stage-relevant feedback observations for evaluated attempts,
including applicable missed tracing portions, incomplete lines or shapes and
incorrect endpoints. & H-AHM-005 & Learning Activity \\
\addlinespace
FR-LRN-016\newline\emph{Adult lesson and starting-stage selection} & Squigglr
shall use the lesson and starting guidance stage selected by the supporting
adult for the learner's session, including an adult-selected later starting
stage. & H-AHM-007 & Learning Activity \\
\addlinespace
FR-LRN-017\newline\emph{Adult-selected practice parameters} & Squigglr shall
apply the learner's adult-selected values for \texttt{STAGE\_\allowbreak PASS},
\texttt{DIFFICULTY\_\allowbreak STREAK}, \texttt{ATTEMPT\_\allowbreak LIMIT},
\texttt{IMPROVE\_\allowbreak WINDOW}, \texttt{ATTEMPT\_\allowbreak EXTENSION}
and \texttt{ATTEMPT\_\allowbreak LIMIT\_\allowbreak MAX}, using the documented
midpoint default for any unset value and rejecting out-of-range or inconsistent
selections. & H-AHM-007 & Learning Activity \\
\addlinespace
FR-LRN-018\newline\emph{Adult authorization for discretionary changes} &
Squigglr shall require explicit authorization from the supporting adult before
changing to another lesson, skipping a guidance stage or moving backward to an
earlier stage; routine pass-based advancement to the next stage of the current
lesson requires no additional adult action. & H-AHM-006 & Learning Activity \\
\addlinespace
FR-LRN-019\newline\emph{Difficulty suggestion} & Squigglr shall offer, without
applying, an optional suggestion to return to an earlier drawing stage when the
learner reaches \texttt{DIFFICULTY\_\allowbreak STREAK} consecutive non-passing
evaluable attempts at an eligible stage and the improvement rule does not
suppress suggestions. When an adult declines or ignores the suggestion, another
suggestion shall require a new full failure streak. & H-AHM-006, H-AHM-007 &
Learning Activity \\
\addlinespace
FR-LRN-020\newline\emph{Improvement-\allowbreak{}based extension} & Squigglr
shall suppress a difficulty suggestion while the learner qualifies as improving,
and, when the unsuccessful-attempt count first reaches
\texttt{ATTEMPT\_\allowbreak LIMIT}, grant up to \texttt{ATTEMPT\_\allowbreak
EXTENSION} additional evaluable stage attempts if the average result of the most
recent \texttt{IMPROVE\_\allowbreak WINDOW} attempts exceeds the average of the
preceding equally sized window by at least \texttt{IMPROVE\_\allowbreak MARGIN},
subject to \texttt{ATTEMPT\_\allowbreak LIMIT\_\allowbreak MAX}. The extension
shall be granted at most once per stage and expire after the allotted additional
evaluable attempts, even if results continue to improve. & H-AHM-006 & Learning
Activity \\
\addlinespace
FR-LRN-021\newline\emph{Stage attempt limit and lesson end} & Squigglr shall end
the active lesson when its unsuccessful-attempt count reaches
\texttt{ATTEMPT\_\allowbreak LIMIT} without a qualifying extension, or a granted
extension is exhausted, or the allowed total-attempt count reaches
\texttt{ATTEMPT\_\allowbreak LIMIT\_\allowbreak MAX}, and make the furthest
stage reached available for the lesson record. & H-AHM-007 & Learning Activity
\\
\addlinespace
FR-LRN-023\newline\emph{Automatic within-lesson advancement} & Squigglr shall
automatically advance from full-outline tracing to endpoint-only guidance, and
from endpoint-only guidance to independent drawing, upon a passing attempt at
the active stage; a passing independent drawing attempt completes the current
lesson without starting a different lesson. & H-AHM-004 & Learning Activity \\
\addlinespace
FR-CRT-007\newline\emph{Undo a change} & Squigglr shall let the child undo the
changes to the open creative drawing one at a time, most recent first, back to
the drawing as it was when the child opened it. & H-MUH-001 & Creative Drawing
\\
\addlinespace
FR-CRT-008\newline\emph{Redo an undone change} & Squigglr shall let the child
redo undone changes one at a time until the child makes a new change. &
H-MUH-001 & Creative Drawing \\
\addlinespace
FR-CRT-010\newline\emph{Refine a finished drawing} & Squigglr shall refine the
open creative drawing automatically each time the child finishes it while
refinement is turned on, and at no other time. & H-MUH-002 & Creative Drawing \\
\addlinespace
FR-CRT-011\newline\emph{Change only stroke paths} & Squigglr shall change only
the path of each existing stroke when it refines a creative drawing. & H-MUH-002
& Creative Drawing \\
\addlinespace
FR-CRT-012\newline\emph{Stay within the refinement tolerance} & Squigglr shall
keep every refined stroke within \texttt{REFINE\_\allowbreak TOLERANCE} of the
corresponding original stroke. & H-MUH-002, H-MUH-006 & Creative Drawing \\
\addlinespace
FR-CRT-013\newline\emph{Keep the refined drawing editable} & Squigglr shall let
the child edit each stroke of a refined drawing with the same actions, and with
the same results, as a stroke the child drew. & H-MUH-001 & Creative Drawing \\
\addlinespace
FR-CRT-014\newline\emph{Do not refine a stroke twice} & Squigglr shall not
refine again a stroke that an earlier refinement changed, unless the child has
edited that stroke since. & H-MUH-002 & Creative Drawing \\
\addlinespace
FR-CRT-015\newline\emph{Leave unsupported drawings unchanged} & Squigglr shall
leave a finished creative drawing unchanged if it has no strokes or is not a
drawing of an item in the supported set. & H-MUH-002, H-MUH-006 & Creative
Drawing \\
\addlinespace
FR-CRT-016\newline\emph{Show whether a drawing is refined} & Squigglr shall
show, while a creative drawing is open, whether it holds a refinement that has
not been undone or restored. & H-MUH-003, H-MUH-006 & Creative Drawing \\
\addlinespace
FR-CRT-017\newline\emph{Restore the exact original} & Squigglr shall return a
creative drawing, on restore, exactly to the original of its most recent
refinement that has not been undone or restored. & H-MUH-001, H-MUH-003 &
Creative Drawing \\
\addlinespace
FR-CRT-018\newline\emph{Child restores the original} & Squigglr shall let the
child restore the original of the open creative drawing. & H-MUH-001, H-MUH-006
& Creative Drawing \\
\addlinespace
FR-CRT-019\newline\emph{Adult restores the original} & Squigglr shall let the
supporting adult, after passing the adult check of SAF-RAY-001, restore the
original of any creative drawing saved under a learner profile in that adult's
account. & H-MUH-001 & Creative Drawing \\
\addlinespace
FR-CRT-020\newline\emph{Leave the drawing unchanged when refinement fails} &
Squigglr shall make no change to a creative drawing when its refinement fails or
the refinement result violates FR-CRT-011 or FR-CRT-012. & H-MUH-002 & Creative
Drawing \\
\addlinespace
FR-CRT-021\newline\emph{Discard an out-of-date result} & Squigglr shall discard
a refinement result if the drawing has been changed or closed since that
refinement started. & H-MUH-002 & Creative Drawing \\
\addlinespace
FR-CRT-023\newline\emph{Save to the active learner} & Squigglr shall save each
creative drawing, with the originals of its refinements, to the learner profile
that was active when the child drew it, without any action from the child. &
H-MUH-003, H-MUH-004, H-MUH-005 & Creative Drawing \\
\addlinespace
FR-CRT-024\newline\emph{Reopen a drawing as it was left} & Squigglr shall reopen
a saved creative drawing exactly as it was last saved, with every stroke
editable and every refinement that has not been undone or restored still
restorable. & H-MUH-001, H-MUH-005 & Creative Drawing \\
\addlinespace
FR-CRT-025\newline\emph{Keep changes through a lost connection} & Squigglr shall
save, once the connection returns, every change the child made to an open
creative drawing while the connection was lost. & H-MUH-005 & Creative Drawing
\\
\addlinespace
FR-CRT-026\newline\emph{Show only the active learner's drawings} & Squigglr
shall show in the creative space only the creative drawings saved under the
active learner profile. & H-MUH-004 & Creative Drawing \\
\addlinespace
FR-CRT-027\newline\emph{No refinement in guided lessons} & Squigglr shall not
refine any stroke drawn in a guided lesson. & H-MUH-002 & Creative Drawing \\
\addlinespace
FR-CRT-028\newline\emph{Creative drawings are not attempts} & Squigglr shall not
count a creative drawing as an attempt. & H-MUH-003 & Creative Drawing \\
\addlinespace
FR-CRT-029\newline\emph{Show unsaved changes} & Squigglr shall show, while a
creative drawing is open, whether every change to it has been saved. & H-MUH-005
& Creative Drawing \\
\addlinespace
FR-CRT-030\newline\emph{Turn refinement off} & Squigglr shall let the child, and
the supporting adult, turn automatic refinement off and on for a learner
profile. & H-MUH-006 & Creative Drawing \\
\addlinespace
FR-REC-001\newline\emph{Attempt record} & Squigglr shall keep, for each attempt
that ends with an outcome, an attempt record holding the lesson, the stage, the
attempt's result, the outcome (whether the attempt passed its stage, FR-LRN-013)
and the date and time the attempt ended. & H-MUH-003 & Attempt Record \\
\addlinespace
FR-REC-002\newline\emph{Actual strokes} & Squigglr shall keep in each attempt
record the strokes the child drew in the attempt, in the order drawn, as they
were captured. & H-MUH-003 & Attempt Record \\
\addlinespace
FR-REC-003\newline\emph{Difficulties and corrections} & Squigglr shall keep in
each attempt record every predefined difficulty observed during the attempt,
each marked as corrected or not corrected within that attempt. & H-MUH-007 &
Attempt Record \\
\addlinespace
FR-REC-004\newline\emph{Adult-set values in use} & Squigglr shall keep in each
attempt record the value of each adult-set value that was in use for the learner
when the attempt was made. & H-MUH-007 & Attempt Record \\
\addlinespace
FR-REC-005\newline\emph{Interrupted attempt} & Squigglr shall keep an attempt
that an adult stops in Squigglr, by ending the session or switching learner,
after at least one stroke but before its outcome is decided, as an attempt
record marked as interrupted, holding the strokes and predefined difficulties up
to that point and no outcome. & H-MUH-004 & Attempt Record \\
\addlinespace
FR-REC-006\newline\emph{Lesson record} & Squigglr shall keep, each time a
learner starts a lesson, a lesson record holding the lesson, the starting stage,
the furthest stage reached, the number of successful attempts at each stage and
how the lesson ended. & H-MUH-007 & Attempt Record \\
\addlinespace
FR-REC-007\newline\emph{Unfinished lesson} & Squigglr shall keep the lesson
record of a lesson that stops before its last stage is passed and before it
reaches an attempt limit, marked as unfinished, with the furthest stage reached
and the successful attempts at each stage up to that point. & H-MUH-007 &
Attempt Record \\
\addlinespace
FR-REC-008\newline\emph{Each attempt saved once} & Squigglr shall save each
attempt record exactly once, under the lesson record of the lesson in which the
attempt was made. & H-MUH-004 & Attempt Record \\
\addlinespace
FR-REC-009\newline\emph{Saved records unchanged} & Squigglr shall keep the
contents of every saved attempt record, and of every lesson record whose lesson
has ended, unchanged except by deletion. & H-MUH-003, H-MUH-004, H-MUH-007 &
Attempt Record \\
\addlinespace
FR-PRG-001\newline\emph{One learner per summary} & Squigglr shall show in each
progress summary only the records of the one learner profile that the summary
names. & H-MUH-004 & Progress Summary \\
\addlinespace
FR-PRG-002\newline\emph{Progress over time by lesson} & Squigglr shall show, for
each lesson in a progress summary, the learner's lesson records in the order
they were started, each with its starting stage, furthest stage reached,
successful attempts at each stage and how the lesson ended. & H-MUH-008 &
Progress Summary \\
\addlinespace
FR-PRG-003\newline\emph{With and without a guide kept apart} & Squigglr shall
show attempts made without a guide separately from attempts made with a guide in
every part of a progress summary that counts or lists attempts. & H-MUH-003 &
Progress Summary \\
\addlinespace
FR-PRG-004\newline\emph{Recurring difficulties} & Squigglr shall list in each
progress summary every predefined difficulty observed in at least
\texttt{RECUR\_\allowbreak MIN} of the learner's attempts at the same stage of
the same lesson, with the number of those attempts and the number in which the
child corrected it. & H-MUH-008 & Progress Summary \\
\addlinespace
FR-PRG-005\newline\emph{Open the attempts behind the summary} & Squigglr shall
let the adult open, from each figure or item in a progress summary, the attempt
records from which it was computed. & H-MUH-008 & Progress Summary \\
\addlinespace
FR-PRG-006\newline\emph{Attempt record view} & Squigglr shall show, for each
attempt record the adult opens, every value kept in it under FR-REC-001 to
FR-REC-005, with the strokes drawn as recorded. & H-MUH-008 & Progress Summary
\\
\addlinespace
FR-PRG-008\newline\emph{Empty history} & Squigglr shall show, when the learner
has no saved attempt record in the chosen lessons and dates, a progress summary
that states that no practice is recorded and shows no stage, count or
difficulty. & H-MUH-008 & Progress Summary \\
\addlinespace
FR-PRG-009\newline\emph{Interrupted attempts shown apart} & Squigglr shall show
interrupted attempts in a progress summary separately from successful and
unsuccessful attempts. & H-MUH-008 & Progress Summary \\
\addlinespace
FR-PRG-010\newline\emph{Stages not attempted} & Squigglr shall mark each stage
of a lesson that the learner has not attempted, including a stage before an
adult-chosen starting stage, as not attempted, distinct from a stage attempted
without success. & H-MUH-003, H-MUH-008 & Progress Summary \\
\addlinespace
FR-PRG-011\newline\emph{Latest attempt included} & Squigglr shall state on each
progress summary the date and time of the most recent attempt it includes. &
H-MUH-008 & Progress Summary \\
\addlinespace
\midrule
\multicolumn{4}{@{}l}{\textbf{Safety-critical requirements}} \\*
\addlinespace
SAF-RAY-001 & Squigglr shall allow lesson choice, starting-stage choice, learner
switching and acceptance of a go-back suggestion only after an adult check that
a child cannot complete by tapping icons, and shall keep a pending suggestion,
without blocking practice, until an adult accepts or declines it. & H-RAY-001,
H-RAY-006 & Safety-Critical \\
\addlinespace
SAF-RAY-002 & Before a lesson starts, Squigglr shall check the device and
browser against \texttt{SUPPORTED\_DEVICES} (TBD) and, if they do not match,
show an adult-readable notice that results may be unreliable. & H-RAY-002 &
Safety-Critical \\
\addlinespace
SAF-RAY-003 & Squigglr shall keep each completed attempt on the device until the
server confirms it is saved, and shall tell the adult which attempts could not
be saved. & H-RAY-002 & Safety-Critical \\
\addlinespace
SAF-RAY-004 & No record, summary or export shall show a diagnostic label, a
judgement of the child's abilities or a comparison with other children, and
every summary and export shall state that it describes practice only and is not
an assessment or diagnosis. & H-RAY-003 & Safety-Critical \\
\addlinespace
SAF-RAY-005 & Squigglr shall end a lesson at \texttt{ATTEMPT\_LIMIT}, or at
\texttt{ATTEMPT\_LIMIT\_MAX} when extended, without any action from the child,
and shall keep the progress saved up to that point. & H-RAY-004 &
Safety-Critical \\
\addlinespace
SAF-RAY-006 & Squigglr shall record each lesson that ends at
\texttt{ATTEMPT\_LIMIT} or \texttt{ATTEMPT\_LIMIT\_MAX}, and each improvement
extension, show the adult how many times these happened for each lesson, and
keep earlier attempts when a lesson is restarted. & H-RAY-005 & Safety-Critical
\\
\addlinespace
SAF-RAY-007 & Squigglr shall identify the active learner on every practice
screen in a form a child who cannot read can recognize, and shall save every
attempt and suggestion to that learner only. & H-RAY-006 & Safety-Critical \\
\addlinespace
SAF-RAY-008 & Squigglr shall save a child's adult-set values (the SRS table
``Central symbolic constants register'') only within their team-set ranges, with
\texttt{ATTEMPT\_LIMIT} larger than \texttt{DIFFICULTY\_STREAK} and
\texttt{ATTEMPT\_LIMIT\_MAX} at least \texttt{ATTEMPT\_LIMIT}; shall use the
middle of each range when the adult does not choose; shall allow changes only
after the adult check; and shall show the values in use alongside the child's
records. & H-RAY-007 & Safety-Critical \\
\addlinespace
SAF-AHM-001 & Squigglr shall prevent duplicate reports of the same stylus input
event from being counted more than once during evaluation or recording of an
attempt, while preserving legitimate separately drawn strokes, including
overlapping strokes. & H-AHM-002 & Safety-Critical \\
\addlinespace
SAF-AHM-002 & Squigglr shall preserve the relative positions and order of stylus
points within the active drawing surface when the canvas size or supported
device orientation changes, and shall reject input whose coordinates cannot be
mapped reliably to that surface. & H-AHM-003 & Safety-Critical \\
\end{longtable}
}

In total, 65 requirements are HA-derived: 14 Learning Activity, 22
Creative Drawing, nine Attempt Record and ten Progress Summary
requirements, and ten safety-critical requirements. Every row of the three
FMEA tables cites at least one of them. The Roadmap section that follows
states the phase in which each is planned.

\section{Roadmap}

This section states when each HA-derived requirement of
Table~\ref{tab:safety-reqs} is planned to be met. The capstone timeline has
three phases, each named after the course deliverable that ends it: PH-1,
the proof-of-concept (PoC) demonstration, from Nov.~23, 2026; PH-2,
Revision~0 (Rev~0), from Feb.~1, 2027; and PH-3, the final demonstration,
from Mar.~15, 2027. These identifiers are those of the SRS Introduction and
are distinct from the interview participant codes. Each date is the start
of the course window for that
deliverable in the Development Plan (SRC-002). Work planned after the
capstone has the phase Future. A phase states when a requirement is
planned to be met; no requirement is claimed here to be implemented.
The PoC experiments use team-created drawings and prepared attempts
(SRC-002), and the PoC does not need a complete application, so the
requirements planned for PH-2 protect the people who use the product as a
whole.

A functional requirement's phase is the one in its Phase field in the SRS.
The safety-critical requirements have no Phase field, so for each of them
Table~\ref{tab:roadmap} gives the phase of the functions it protects and
the reason. The table gives one reason for each group of requirements that
share a purpose, and names the hazards each group mitigates.

{\footnotesize
\setlength{\tabcolsep}{4pt}
\exhyphenpenalty=10000 % keep IDs such as FR-CRT-017 on one line
\begin{longtable}{@{}
  >{\raggedright\arraybackslash}p{1.7cm}
  >{\raggedright\arraybackslash}p{2.1cm}
  >{\raggedright\arraybackslash}p{\dimexpr\linewidth-3.8cm-4\tabcolsep\relax}
  @{}}
\caption{Phase of each HA-derived requirement}
\label{tab:roadmap}\\
\toprule
\textbf{Phase (date)} & \textbf{Requirement IDs} & \textbf{Rationale} \\
\midrule
\endfirsthead
\toprule
\textbf{Phase (date)} & \textbf{Requirement IDs} & \textbf{Rationale} \\
\midrule
\endhead
\midrule
\multicolumn{3}{@{}r@{}}{\emph{Continued on the next page}} \\
\endfoot
\bottomrule
\endlastfoot
PH-1 (Nov.~23, 2026) & FR-CRT-007, FR-CRT-010, FR-CRT-011, FR-CRT-012,
FR-CRT-013, FR-CRT-017, FR-CRT-018 &
Editable refinement that can be undone is the main technical uncertainty
of the PoC. Its second experiment refines drawings of the agreed supported
set into editable paths, modifies the result, undoes the refinement and
compares the restored drawing with the saved original (SRC-002). These
requirements state what that experiment must show, so the controls on a
creative change that cannot be undone or is applied wrongly are planned
from the first phase (H-MUH-001, H-MUH-002, H-MUH-003, H-MUH-006). \\
\addlinespace
PH-1 (Nov.~23, 2026) & FR-CRT-020 &
The third PoC experiment must show that an invalid instruction from a
candidate model is rejected and leaves the drawing unchanged (SRC-002).
This requirement states that behaviour. If the model investigation is
deferred, as the Development Plan allows, the requirement still applies
to a failed geometric refinement (H-MUH-002). \\
\addlinespace
PH-1 (Nov.~23, 2026) & FR-REC-001, FR-REC-003 &
The fourth PoC experiment inspects the attempt records to confirm that the
observations and the assistance used remain available (SRC-002). These two
requirements define that record: each attempt's lesson, stage and outcome,
and every predefined difficulty observed, marked as corrected or not
(H-MUH-003, H-MUH-007). \\
\addlinespace
PH-1 (Nov.~23, 2026) & FR-LRN-007, FR-LRN-010, SAF-AHM-001, SAF-AHM-002 &
The first PoC experiment draws with touch and a compatible stylus on the
target tablet, repeats the interactions on an additional device and browser,
and checks that strokes follow the input position and remain continuous
(SRC-002). These requirements state what the capture must keep: the order
and boundaries of strokes (FR-LRN-007), no passing status for missing or
ambiguous input (FR-LRN-010), no repeated count of one input event
(SAF-AHM-001) and no distortion of stroke positions when the canvas size or
device orientation changes (SAF-AHM-002) (H-AHM-001, H-AHM-002,
H-AHM-003). \\
\addlinespace
PH-1 (Nov.~23, 2026) & FR-LRN-011, FR-LRN-013, FR-LRN-016, FR-LRN-023 &
The fourth PoC experiment uses prepared sequences of successful and
unsuccessful attempts and demonstrates reduced assistance following success
(SRC-002). These requirements state when an attempt is complete and passes
(FR-LRN-011, FR-LRN-013), that the lesson and starting stage the adult chose
are used (FR-LRN-016) and that a passing attempt advances within the lesson
without a further adult action (FR-LRN-023) (H-AHM-004, H-AHM-005,
H-AHM-007). \\
\midrule
PH-2 (Feb.~1, 2027) & FR-CRT-014, FR-CRT-015, FR-CRT-016, FR-CRT-019,
FR-CRT-021 &
Rev~0 demonstrates the product as a whole, after the focused PoC
experiments. These requirements bound refinement
in the cases the PoC experiments do not cover: a stroke that would be
refined twice, a drawing outside the supported set, a result that arrives
after the drawing has changed, a refinement that cannot be seen, and a
restore by the adult after the adult check of SAF-RAY-001 (H-MUH-001,
H-MUH-002, H-MUH-003, H-MUH-006). \\
\addlinespace
PH-2 (Feb.~1, 2027) & FR-CRT-023, FR-CRT-024, FR-CRT-026, FR-CRT-027,
FR-CRT-028 &
Saving creative drawings to learner profiles, and keeping them apart from
guided lessons and attempts, needs adult accounts, learner profiles and
stored records. The PoC does not need a complete application (SRC-002), so
these requirements are planned for Rev~0 (H-MUH-001, H-MUH-002, H-MUH-003,
H-MUH-004, H-MUH-005). \\
\addlinespace
PH-2 (Feb.~1, 2027) & FR-CRT-030 &
Refinement is turned off for a learner profile, so this requirement needs
the learner profiles and settings planned for Rev~0, like the requirements
that keep creative drawings to one learner. Until then, every refinement
stays within \texttt{REFINE\_TOLERANCE} and can be undone or restored
(FR-CRT-007, FR-CRT-012, FR-CRT-017, FR-CRT-018, PH-1) (H-MUH-006). \\
\addlinespace
PH-2 (Feb.~1, 2027) & FR-REC-002, FR-REC-004 to FR-REC-009 &
Complete attempt and lesson records, each saved once and unchanged after
saving, are needed as soon as adults rely on the records (G.5). The values
kept by FR-REC-004 have ranges that the SRS table ``Central symbolic
constants register'' defines; their verification is due Nov.~2, 2026
(OI-05) (H-MUH-003, H-MUH-004, H-MUH-007). \\
\addlinespace
PH-2 (Feb.~1, 2027) & FR-PRG-001 to FR-PRG-006, FR-PRG-008 to FR-PRG-011 &
Progress summaries are built from the complete records, so they are
planned with them (G.6). These requirements keep each summary to one
learner and consistent with the records behind it (H-MUH-003, H-MUH-004,
H-MUH-008). \\
\addlinespace
PH-2 (Feb.~1, 2027) & SAF-RAY-001, SAF-RAY-007 &
Adult-only actions, learner switching and an active learner exist only
once the product has adult accounts and several learner profiles, which
the PoC does not need. Both are planned with the requirements that keep
drawings and summaries to one learner (FR-CRT-023, FR-CRT-026,
FR-PRG-001) (H-RAY-001, H-RAY-006). \\
\addlinespace
PH-2 (Feb.~1, 2027) & SAF-RAY-002 &
The device check needs the list of supported devices and browsers
(\texttt{SUPPORTED\_DEVICES}), which the team sets by Nov.~2, 2026
(OI-09). The first PoC experiment records the hardware and browser
versions tested (SRC-002), which informs that list (H-RAY-002). \\
\addlinespace
PH-2 (Feb.~1, 2027) & SAF-RAY-003 &
Keeping each completed attempt on the device until the server confirms
the save protects the stored records, which are planned for Rev~0. The
requirement is planned with FR-REC-008 and FR-PRG-011 (H-RAY-002). \\
\addlinespace
PH-2 (Feb.~1, 2027) & SAF-RAY-004 &
The requirement applies to every record view and progress summary, which
are planned for Rev~0 with the Progress Summary requirements (H-RAY-003). \\
\addlinespace
PH-2 (Feb.~1, 2027) & SAF-RAY-005, SAF-RAY-006 &
Ending a lesson at its limit, and recording lessons that reach the limit
and are restarted, belong to the complete lesson flow and the lesson
record. The fourth PoC experiment, as the Development Plan describes it,
demonstrates changes in assistance rather than lesson ends (SRC-002).
Both are planned with FR-REC-006 and FR-REC-007, and SAF-RAY-005 with
FR-LRN-021, which ends the lesson at its attempt limit (H-RAY-004,
H-RAY-005). \\
\addlinespace
PH-2 (Feb.~1, 2027) & SAF-RAY-008 &
Checking adult-set values uses their team-set ranges, which the SRS table
``Central symbolic constants register'' defines and whose verification is
due Nov.~2, 2026 (OI-05). Applying the adult's values is planned with
FR-LRN-017, and showing the values in use with the records with FR-REC-004
(H-RAY-007). \\
\addlinespace
PH-2 (Feb.~1, 2027) & FR-LRN-008, FR-LRN-009 &
Attempts of several strokes, and attempts that are interrupted or
unfinished, extend the stroke capture planned for PH-1 (FR-LRN-007). Their
phase is the one in their Phase field in the SRS (H-AHM-001, H-AHM-002,
H-AHM-006). \\
\addlinespace
PH-2 (Feb.~1, 2027) & FR-LRN-014, FR-LRN-017 to FR-LRN-021 &
Feedback observations, the adult's chosen values, adult authorization of
discretionary changes, the difficulty suggestion, the improvement extension
and the end of a lesson at its attempt limit belong to the complete lesson
flow. Their phase is the one in their Phase field in the SRS, and it matches
that of SAF-RAY-005, SAF-RAY-006 and SAF-RAY-008, which protect the same
functions (H-AHM-005, H-AHM-006, H-AHM-007). \\
\midrule
PH-3 (Mar.~15, 2027) & FR-CRT-008 &
Redo (priority Could) only recovers a change undone by mistake. Until PH-3,
every change can still be undone (FR-CRT-007, PH-1), and the original of a
refinement can be restored (FR-CRT-017, FR-CRT-018, PH-1) (H-MUH-001). \\
\addlinespace
PH-3 (Mar.~15, 2027) & FR-CRT-025 &
Keeping the changes made while the connection is lost (priority Should)
is the only requirement that saves them (H-MUH-005). Until PH-3, every
creative drawing is saved without any action from the child and reopens
as it was left (FR-CRT-023, FR-CRT-024, PH-2), so only the changes made
while the connection is lost are at risk. Offline use is out of scope
(ASM-005), and the FMEA's test of closing the page during a lost
connection records which changes are kept (H-MUH-005). \\
\addlinespace
PH-3 (Mar.~15, 2027) & FR-CRT-029 &
Showing whether every change has been saved (priority Should) lets an adult
keep the page open until a change made while the connection was lost has
been saved, because offline use is out of scope (ASM-005) and such changes
are kept only while the page stays open (FR-CRT-025) (H-MUH-005). \\
\end{longtable}
}

No HA-derived requirement is planned for Future: every row of the three
FMEA tables has its mitigating requirements planned within the capstone.
Four of
the items left outside the scope of this analysis, as the Scope and
Purpose of Hazard Analysis section explains, are features that may be
added later. Each would need hazards and safety
requirements of its own before it is built, so those safety requirements
are Future work:
\begin{itemize}
  \item \emph{Camera-based finger tracing}, a stretch goal (SRC-001): a
  camera would capture images of the child, while a learner profile holds
  only a first name and a non-identifying username (BR-001), so its
  privacy hazards must be analysed first.
  \item \emph{Letter tracing and drawing}, a stretch goal (SRC-001): the
  existing requirements apply to letter lessons, and letter-specific
  hazards, such as an adult reading letter attempts as an assessment
  (BR-006), need requirements when letters are specified.
  \item \emph{Finger input as a supported input} (CON-003, LC-006): stage
  checks and value ranges are set for stylus input, so the canvas and
  lesson hazards must be reviewed before finger input is supported.
  \item \emph{An offline mode} (ASM-005): changes and attempts kept
  without a connection would extend FR-CRT-025 and SAF-RAY-003, whose
  hazards are H-MUH-005 and H-RAY-002.
\end{itemize}

\newpage{}

\section{Appendix --- Generative AI Use Disclosure}

\wss{ChatGPT (version 5) was used to brainstorm ideas for the template for this disclose section.}

\subsection{AI Tools Used}

\wss{Give the name of the AI tool(s) used and the version.}

\subsection{How and Where Used}

\wss{Explain what the tool(s) were used for.  Examples include: brainstorming,
explaining a technical concept, reviewing the document, generating or modifying
text, generating or modifying diagrams, identifying errors or omissions or
inconsistencies.}

\wss{For each of the uses, identify which parts of the document were affected.}

\wss{You do not need to report every interaction, but you should disclose the 
uses that materially contributed to your work}

\subsection{Verification of AI-Generated Content}

\wss{\begin{itemize}
    \item Describe how the team checked AI-generated or AI-assisted material
    before including it in the document.
    \item In particular, verify factual claims, technical assertions,
    requirements, calculations, references, and other information for which an
    AI system may produce incorrect or fabricated results.
    \item \textbf{The team is responsible for the accuracy, completeness,
consistency, and appropriateness of the submitted document, regardless of
whether material was generated by a human or an AI system.}
\end{itemize}}

\subsection*{Individual Disclosures}

\textbf{Rayan:} I used Anthropic's Claude Opus 5.5 at medium effort,
accessed through Claude Code, while preparing my sections of this document.
I developed the ideas and made the project decisions in my sections myself,
and Moustafa and I conducted the specialist interview and follow-up
questionnaire. A major use of AI was organizing our sources: I used the
model to sort through the interview transcript, questionnaire answers, and
meeting notes, and to extract the specific statements, with their
timestamps, that supported the ideas I wanted to include. This would have
been very difficult to do by hand given the length of the material. I also
used it to turn my notes and decisions into draft wording, to improve
sentence structure, clarity, and flow, and to help with \LaTeX{}
formatting. Through Claude Code, it also helped manage the repository
(branches, commits, and merges) and helped me review my teammates' pull
requests by checking their changes against our decisions and the rest of
the document; I decided which review comments to post.

All ideas and decisions in my sections are my own. I reviewed every
AI-assisted passage against my intended meaning and checked each cited
statement against the original transcript, questionnaire, or meeting notes
to confirm that it says what we claim it says.

\newpage{}

\section{Appendix --- Reflection}

\wss{Not required for CAS 741}

\input{../Reflection.text}

\begin{enumerate}

\item What went well while writing this deliverable?

\textbf{Rayan:} Identifying the hazards was easier than expected because we
already had good sources to work from. The transcript of our interview with
three domain specialists on October 1, 2026 gave us concrete statements, with
timestamps, about how the intended children use apps and what could go wrong
for them. For example, the specialists said adults should decide progression
rather than the app, that audio feedback should be possible to mute, and that
Squigglr may be used at home, at school or in other settings. Each of these
pointed directly to a failure mode we could analyse, so we were not guessing at
hazards from scratch. Our meeting with our supervisor, Dr.~Anand, on October
14, 2026 helped in the same way: once Dr.~Anand explained that thresholds
should be set by the adult for each child within a range, the hazards around
wrong or extreme threshold values became clear. Because the SRS already
recorded these sources, business rules and assumptions, we could trace each
hazard and critical assumption back to a specific statement, which also made
the mitigations easier to justify.

\item What pain points did you experience during this deliverable, and how
did you resolve them?

\textbf{Rayan:} The main pain point was that my hazards kept changing as other
parts of the documents changed. I wrote my first set of hazards and critical
assumptions based on the SRS as it stood at the time, where the team would set
fixed values for thresholds such as \texttt{DIFFICULTY\_STREAK} and
\texttt{ATTEMPT\_LIMIT}. In our meeting on October 14, 2026, Dr.~Anand advised
that suitable values differ from child to child, so I changed those thresholds
to be set by the supporting adult within a range the team defines. When I
reviewed the hazard analysis again, I found that this changed the hazards: the
risk was no longer only a wrong team-chosen value, but also an adult choosing a
value at the edge of the range, so I added a new hazard and safety requirement
for it. The same thing happened when I added the improvement rule, which can
extend a lesson while a child is improving, and when we confirmed that the
child does not navigate inside a lesson. Each time, I had to go back and update
the affected hazards, critical assumptions and mitigations so they matched the
SRS. It was time-consuming, but it showed me how closely the hazard analysis
depends on the requirements, and that it has to be reviewed again whenever the
SRS changes.

\item Which of your listed risks had your team thought of before this
deliverable, and which did you think of while doing this deliverable? For
the latter ones (ones you thought of while doing the Hazard Analysis), how
did they come about?

\textbf{Rayan:} Before this deliverable, the team had already thought of two
of my risks. Because Squigglr runs in a browser, we expected that attempts
could be lost or not saved on an unreliable browser or device (H-RAY-002). We
also knew that adults might read a child's records as a diagnosis (H-RAY-003),
which is why the problem statement already avoids any diagnostic claims.

The other five risks I only thought of while doing the hazard analysis. They
came about because this was the first time I had to work out the requirements
for the application in detail, while writing the SRS alongside this document.
Most of them came from the critical assumptions, and several of those
assumptions came from our interview with the specialists on October 1, 2026.
The specialists said children can use the app through images and icons while
adults handle setup, which led to the assumption that an adult sets up each
session. Asking what happens if no adult is actually watching led to the risk
of a child reaching adult-only actions (H-RAY-001), and, together with defining
what happens after repeated failed attempts, to the risks of a child becoming
distressed (H-RAY-004) and of a lesson being restarted straight after it ends
(H-RAY-005). The specialists also said Squigglr may be used at home, at school
or in other settings, which led to the assumption that a device may be shared
and to the risk of a child's records being saved under the wrong profile
(H-RAY-006). The last risk came from our meeting on October 14, 2026, where
Dr.~Anand advised that adults should set the thresholds for each child; this
raised the risk of a threshold being set to a value that does not suit the
child (H-RAY-007).

Looking back, most of my risks only became visible once the requirements were
concrete enough to question. The risks we knew about beforehand were general
ones that apply to almost any browser app, while the ones I found during this
deliverable were specific to how children and adults will actually use
Squigglr. This showed me that the hazard analysis cannot be done well in
isolation; it depends on knowing the requirements and on hearing from people
who understand the users.

\item Other than the risk of physical harm (some projects may not have any
appreciable risks of this form), list at least 2 other types of risk in
software products. Why are they important to consider?

\textbf{Rayan:} Most software carries risks that have nothing to do with
physical harm. One is privacy and security. Almost every application stores
some information about its users, and if that information is exposed, accessed
by the wrong person or mixed up between users, it can harm those users and
destroy their trust in the product. This matters even more when the users are
vulnerable. In Squigglr, for example, the records describe a child's
difficulties, so they must only be visible to the adults responsible for that
child. A second is data integrity and availability. People make decisions based
on what software shows them, so data that is lost, corrupted or shown
incorrectly can lead them to the wrong decision, even when nothing visibly
crashes. In our case, an adult choosing a child's next lesson relies on the
attempt records being complete and accurate. A third is psychological harm.
Software that frustrates, confuses or misleads its users can affect their
well-being or lead them to wrong conclusions, for example by implying something
about a person that the software cannot actually judge. For Squigglr, that
means avoiding repeated failure without a break and making sure progress
summaries are never read as a diagnosis. These risks are important to consider
because they are easy to overlook when there is no obvious physical danger, yet
they can cause real harm to users and to the people who rely on the software.

\end{enumerate}

\end{document}